% Exercise info
\newcommand{\coursename}{Compiler Construction}
\newcommand{\semesterlong}{Spring Semester 2026}
\newcommand{\semestershort}{Spring 2026}
\newcommand{\exercisedue}{Tuesday, March 17}
\newcommand{\exercisedate}{Tuesday, March 10}
\newcommand{\exercisenumber}{03}
\newcommand{\exercisetitle}{A Generated Scanner}

\newcommand{\teachingsite}{ILIAS\xspace}
\newcommand{\contact}{sandro.hernandezgoicochea@unibe.ch,timo.kehrer@unibe.ch\xspace}

% Draw header
\begin{center}
	% Logos
% 	\noindent\hspace{1cm}\includegraphics[height=2.5cm]{figures/logo-se}\hfill{}\includegraphics[height=2.5cm]{figures/logo-uni}\hspace*{1cm}%\vspace{0.5cm}
	\hfill{}\includegraphics[height=1.75cm]{../formatting/design/logo-uni}\hspace*{1cm}
	\vspace{6ex}

	
	{\Large\textbf{Exercise \exercisenumber: \exercisetitle}}\\
	
	\vspace{2ex}
	{\large{\coursename\ (\semestershort)}} \\
	\vspace{2ex}
	Handout: \exercisedate\ \\
	Discussion: \exercisedue\ \\
\end{center}

% Instruction
\hrule
%\noindent


% \noindent
% Dieses \"Ubungsblatt beinhaltet \numquestions\ Aufgaben mit einer Gesamtzahl von \numpoints\ Punkten.

\vspace{2ex}
